\documentclass[11pt,a4paper]{article}

\usepackage[margin=1in]{geometry}
\usepackage{amsmath,amssymb,amsthm}
\usepackage{mathtools}
\usepackage[dvipsnames]{xcolor}
\usepackage{hyperref}
\usepackage{booktabs}
\usepackage{array}
\usepackage{enumitem}

\hypersetup{
  colorlinks=true,
  linkcolor=MidnightBlue,
  citecolor=MidnightBlue,
  urlcolor=MidnightBlue
}

\newtheorem{conjecture}{Conjecture}
\newtheorem{problem}{Open Problem}
\theoremstyle{definition}
\newtheorem{definition}[conjecture]{Definition}
\theoremstyle{remark}
\newtheorem{remark}[conjecture]{Remark}

\newcommand{\PCF}{\mathrm{PCF}}
\newcommand{\Q}{\mathbb{Q}}
\newcommand{\Z}{\mathbb{Z}}
\newcommand{\R}{\mathbb{R}}
\newcommand{\C}{\mathbb{C}}
\newcommand{\disc}{\mathrm{disc}}
\newcommand{\Gal}{\mathrm{Gal}}
\newcommand{\Spec}{\mathrm{Spec}}

\title{PCF-2: A Program Statement for the Cubic Extension\\
       of the $\Delta$-Discriminant Transcendence Predicate}
\author{Papanokechi%
\thanks{ORCID: \href{https://orcid.org/0009-0000-6192-8273}{0009-0000-6192-8273}.
Independent Researcher, Yokohama, Japan.
\textbf{AI use disclosure.} This program statement was prepared using the
AEAL/SIARC multi-agent methodology. GitHub Copilot (powered by Anthropic
Claude Opus 4.7) and Anthropic Claude were used to assist with literature
synthesis, drafting of expository prose, and consistency-checking against
prior preprints in the series. All conjecture formulations, scope
decisions, and editorial choices were made by the human author, who takes
full responsibility for the content. The AI tools are not authors. See the
AEAL methodology of~\cite{P8AEAL}.}}
\date{May 1, 2026}

\begin{document}
\maketitle

\begin{abstract}
This paper is a \emph{program statement}, not a results paper. It frames
PCF-2 (P13 in the SIARC stack), the cubic extension of
PCF-1~\cite{Papanokechi2026PCF1} from degree-two polynomial continued
fractions to degree-three. PCF-1 established a sharp empirical dichotomy
across $30$ degree-two families: the sign of the balanced discriminant
$\Delta_{2}=\beta^{2}-4\alpha\gamma$ predicts whether the limit of the
PCF admits an elementary closed form ($\Delta_{2}>0$, $24/24$ Trans-stratum
families) or fails PSLQ detection against an $18$-element basis at $220$
digits ($\Delta_{2}<0$, $6/6$ CM candidates), with the Stokes-exponent
predicate $S<1$ confirmed across all six $\Delta<0$ families. The cubic
case is qualitatively different. The natural invariant of the denominator
polynomial $b(n)=\alpha_{3}n^{3}+\alpha_{2}n^{2}+\alpha_{1}n+\alpha_{0}$ is
the cubic discriminant $\Delta_{3}$, which is no longer a sign invariant
but a quantity whose sign \emph{and} Galois-group structure
($S_{3}$, $A_{3}=C_{3}$, or split) jointly determine the splitting field
and the analytic obstruction. We propose four open conjectures:
(i) a cubic analogue of the sign dichotomy with the Galois-group axis;
(ii) a Galois-stratified CM-field structure replacing the binary
imaginary-quadratic predicate;
(iii) a Stokes signal calibrated for irreducible-cubic-CM families;
and (iv) an alignment with the Painlev\'e IV equation rather than P-III,
motivated by the cubic Galois-group / Painlev\'e hierarchy correspondence.
We sketch a four-to-six-week computational plan that reuses the
infrastructure of PCF-1 and identifies three concrete first sessions
(PCF-2 A/B/C). PCF-2 is positioned to either confirm that the
$\Delta$-predicate of PCF-1 extends to higher degree (universality) or to
identify degree two as a special boundary (specificity).
\end{abstract}

\noindent\textbf{Series label:} P13 / PCF-2.\;
This program statement is P13 in the AI Epistemic Governance Stack and
simultaneously PCF-2 in the mathematical companion track. It anchors a
research program; it does not contain new computational results. The
companion document is the SIARC umbrella program
statement~\cite{Papanokechi2026UMB}.

\tableofcontents

\section{Position in the SIARC Stack}\label{sec:position}

The SIARC research program organizes polynomial continued fractions (PCFs)
into an arithmetic stratification indexed by the dominant growth degree
and an obstruction tier. The program-level architecture is articulated in
the umbrella program statement~\cite{Papanokechi2026UMB}; the
$d=2$ instantiation is PCF-1~\cite{Papanokechi2026PCF1}; the present
document opens the $d=3$ instantiation.

\paragraph{What PCF-1 established.}
For PCFs of the form
\[
K(a_{n},b_{n})=b_{0}+\cfrac{a_{1}}{b_{1}+\cfrac{a_{2}}{b_{2}+\cdots}}
\]
with $a_{n}\in\Z[n]_{\le 1}$ and $b_{n}\in\Z[n]_{\le 2}$, PCF-1
proposed the sign of the balanced discriminant
$\Delta_{2}=\beta^{2}-4\alpha\gamma$ as a transcendence predicate. The
core empirical finding was a sharp dichotomy across $30$ families: $24/24$
Trans-stratum families with $\Delta_{2}\in\{+1,+8\}$ admit elementary
closed forms via Lindemann--Weierstrass, while all $6/6$ CM candidates
with $\Delta_{2}\in\{-3,-4,-7,-11,-20,-28\}$ return PSLQ non-detection
against an $18$-element basis at working precision $220$ digits, and all
six exhibit the Stokes-exponent signature $S<1$ under a CM-respecting
deformation. The V$\_$quad family ($\Delta_{2}=-11$) is the explicit
Painlev\'e III$(D_{6})$ prototype. PCF-1 v1.2 also showed (via QL15,
$\Delta_{2}=-20$, class number $h=2$, non-Heegner; and the intra-field
pair QL06/QL26 in $\Q(\sqrt{-7})$ at $h=1$ versus $h=2$) that the
controlling invariant is the imaginary-quadratic field $\Q(\sqrt{\Delta})$
itself, not its class number~\cite{Papanokechi2026PCF1}.

\paragraph{Why degree three is the natural next step.}
Three structural features make $d=3$ the canonical extension target.

\begin{enumerate}
\item \emph{Discriminant lifts naturally.} The cubic discriminant
$\Delta_{3}$ of $b(n)=\alpha_{3}n^{3}+\alpha_{2}n^{2}+\alpha_{1}n+\alpha_{0}$
is a classical $\mathrm{SL}_{2}(\Z)$-invariant of degree $4$ in the
coefficients (Definition~\ref{def:Delta3}). Its sign distinguishes
totally real cubic splitting from one-real-two-complex splitting, the
direct analogue of the $d=2$ sign dichotomy.
\item \emph{Galois group is no longer trivial.} A degree-three irreducible
polynomial has Galois group $S_{3}$ or $A_{3}=C_{3}$ over $\Q$,
distinguished by whether $\Delta_{3}$ is a square in $\Q$. This adds an
arithmetic axis absent in the degree-two case (where $\Gal(b/\Q)$ is
$C_{2}$ whenever $b$ is irreducible).
\item \emph{Painlev\'e hierarchy aligns with cubic Galois structure.} The
classical Painlev\'e equations admit a cubic-singularity analogue (P-IV)
whose monodromy data is naturally indexed by cubic Galois groups, in
the same way P-III governs the imaginary-quadratic CM case in PCF-1.
\end{enumerate}

\paragraph{What PCF-2 will not address.}
PCF-2 does not extend automatically to degree four or higher. The cubic
case is the unique additional case in which the Galois group is finite
and explicitly classifiable, and in which the Painlev\'e correspondence
is well-posed. Extensions to $d=4$ and beyond are deferred to PCF-3 and
later instantiations of the SIARC program.

\section{Setup: Degree-Three Polynomial CFs}\label{sec:setup}

We work with PCFs of the form
\begin{equation}\label{eq:pcfd3}
K(a_{n},b_{n})=b_{0}+\cfrac{a_{1}}{b_{1}+\cfrac{a_{2}}{b_{2}+\cdots}}
\end{equation}
with
\begin{align}
a_{n} &= \delta_{1}\,n + \delta_{0},\qquad \delta_{1},\delta_{0}\in\Z, \label{eq:an}\\
b_{n} &= \alpha_{3}\,n^{3} + \alpha_{2}\,n^{2} + \alpha_{1}\,n + \alpha_{0},\qquad \alpha_{i}\in\Z,\ \alpha_{3}\ne 0. \label{eq:bn}
\end{align}
The convergents satisfy the Wallis recurrence $p_{n}=b_{n}p_{n-1}+a_{n}p_{n-2}$
with the seeds $(p_{-1},p_{0},q_{-1},q_{0})=(1,b_{0},0,1)$, identical to
the $d=2$ convention of PCF-1 to facilitate code reuse. Coefficient
ordering for the family $b_{n}$ is $[\alpha_{3},\alpha_{2},\alpha_{1},\alpha_{0}]$
(leading first), matching the PCF-1 \texttt{f1\_base\_computation.py}
convention.

\begin{definition}[Cubic balanced discriminant]\label{def:Delta3}
For a balanced PCF in the $d=3$ class with $b_{n}$ as in~\eqref{eq:bn},
\[
\Delta_{3}\;=\;\disc(b)\;=\;
18\alpha_{3}\alpha_{2}\alpha_{1}\alpha_{0}
-4\alpha_{2}^{3}\alpha_{0}
+\alpha_{2}^{2}\alpha_{1}^{2}
-4\alpha_{3}\alpha_{1}^{3}
-27\alpha_{3}^{2}\alpha_{0}^{2}.
\]
We say a degree-three PCF is \emph{$\Z$-primitive} if
$\gcd(\alpha_{3},\alpha_{2},\alpha_{1},\alpha_{0})=1$. We say it is
\emph{irreducible} if $b$ is irreducible over $\Q$, and we always
restrict to the irreducible $\Z$-primitive case in this program.
\end{definition}

\begin{remark}\label{rem:primitive}
The PCF-1 normalization $\gcd(\alpha,\beta,\gamma)=1$ extended to $d=3$
removes the $1$-parameter integer-scaling redundancy and forces
$\Delta_{3}$ to be a $\Z$-valued invariant of the family, not a class
representative. The cubic discriminant is invariant under
$\alpha_{i}\mapsto\alpha_{i}+c\binom{i}{1}$-style $n$-shifts only up to a
predictable scaling; the convention~\eqref{eq:bn} fixes the
$n$-coordinate.
\end{remark}

\paragraph{Splitting types and Galois groups.}
For an irreducible $\Z$-primitive cubic $b\in\Z[n]$:
\begin{itemize}
\item If $\Delta_{3}>0$, then $b$ has three real roots and the splitting
field is the cubic resolvent $\Q(\theta)$ for $\theta$ a real root of $b$,
together with $\sqrt{\Delta_{3}}\in\R$. The Galois group is
$\Gal(b/\Q)=A_{3}=C_{3}$ when $\Delta_{3}$ is a perfect square in $\Q$,
and $\Gal(b/\Q)=S_{3}$ otherwise.
\item If $\Delta_{3}<0$, then $b$ has one real root and a complex
conjugate pair, the splitting field contains $\Q(\sqrt{\Delta_{3}})$
(an imaginary-quadratic subfield), and the Galois group is necessarily
$\Gal(b/\Q)=S_{3}$.
\end{itemize}
This trichotomy --- $A_{3}$ totally real, $S_{3}$ totally real,
$S_{3}$ complex --- is the cubic analogue of the binary
$\Delta_{2}>0$ vs.\ $\Delta_{2}<0$ split, and it is what gives PCF-2
genuinely new content beyond PCF-1.

\paragraph{Spec$(K)$ in the cubic case.}
We retain the v5 schema of PCF-1,
\[
\Spec(K)=(d,\Lambda,\Delta,\rho,\tau_{\mathrm{mono}},\tau_{\mathrm{arith}};S,\sigma),
\]
with $d=3$ now, and with $\Delta=\Delta_{3}$ as defined above. The
Poincar\'e--Perron normalized growth rate $\Lambda$ in the balanced
$d=3$ class satisfies the cubic
\begin{equation}\label{eq:poincare3}
\Lambda^{3}-\alpha_{3}\Lambda^{2}-\alpha_{3}\delta_{1}\Lambda-\alpha_{3}\delta_{0}=0,
\end{equation}
or a closely related cubic depending on the precise scaling
convention; the analytic determination of~\eqref{eq:poincare3} is part
of PCF-2 Session A. The arithmetic sector $\tau_{\mathrm{arith}}$ now
records the splitting field of $b$ together with the Galois group
$\Gal(b/\Q)\in\{C_{3},S_{3}\}$.

\section{Predictions and Open Conjectures}\label{sec:conjectures}

We state four open conjectures, ordered from most direct lift of PCF-1 to
most speculative.

\begin{conjecture}[B$_{3}$ part (i): cubic sign dichotomy]\label{conj:B1}
Let $K$ be a balanced $d=3$ PCF with irreducible $\Z$-primitive
denominator $b$ as in~\eqref{eq:bn}, and let $\Delta_{3}=\disc(b)$.

\emph{(A)} If $\Delta_{3}<0$ (one real root, $\Gal(b/\Q)=S_{3}$), then
$L=K(a_{n},b_{n})$ admits no elementary closed form: equivalently,
$L\notin\Q(\pi,e,\log\Q^{\times},\zeta(3),\sqrt{\Q},\gamma)$ in the sense
of PSLQ non-detection against the PCF-1 $18$-element basis at working
precision $\ge 220$ digits and integer bound $\ge 10^{10}$.

\emph{(B)} If $\Delta_{3}>0$ \emph{and} $\sqrt{\Delta_{3}}\in\Q$ (so
$\Gal(b/\Q)=C_{3}$, the totally-real cyclic case), then
$L$ \emph{may} admit a closed form; this is the cubic analogue of the
PCF-1 $\Delta_{2}>0$ Trans-stratum and is the part of the conjecture
most likely to admit explicit Lindemann--Weierstrass evaluation.

\emph{(C)} If $\Delta_{3}>0$ \emph{and} $\sqrt{\Delta_{3}}\notin\Q$ (so
$\Gal(b/\Q)=S_{3}$, the totally-real non-cyclic case), then $L$ admits no
elementary closed form, with the same PSLQ-non-detection criterion as
(A).
\end{conjecture}

\begin{remark}\label{rem:trichotomy}
Conjecture~\ref{conj:B1} is a \emph{trichotomy} rather than a sign
dichotomy. The new ingredient is the Galois-group axis: the cyclic case
$C_{3}$ is the sole closed-form-admitting cell. This is the structural
prediction PCF-2 Sessions A and B are designed to falsify or confirm.
\end{remark}

\begin{conjecture}[B$_{3}$ part (ii): Galois-stratified CM-field structure]\label{conj:B2}
Let $K$ be a balanced $d=3$ PCF with irreducible $\Z$-primitive
denominator $b$. Let $\theta_{1},\theta_{2},\theta_{3}\in\C$ be the roots
of $b$, and let $L_{b}/\Q$ be the splitting field.
\begin{enumerate}
\item If $\Delta_{3}<0$, then $L_{b}=\Q(\theta,\sqrt{\Delta_{3}})$ where
$\theta$ is the unique real root, $[L_{b}:\Q]=6$, and the imaginary
quadratic subfield $\Q(\sqrt{\Delta_{3}})\subset L_{b}$ controls the
Stokes data.
\item If $\Delta_{3}>0$ and $\Gal(b/\Q)=C_{3}$, then $L_{b}=\Q(\theta_{1})$
is a totally-real cyclic cubic field of degree~$3$.
\item If $\Delta_{3}>0$ and $\Gal(b/\Q)=S_{3}$, then $L_{b}=\Q(\theta_{1},\sqrt{\Delta_{3}})$
is a totally-real non-cyclic sextic field.
\end{enumerate}
The singular points of the associated rank-three Heun-type ODE for $K$
all lie in $L_{b}$.
\end{conjecture}

\begin{remark}\label{rem:not-binary}
Unlike the PCF-1 case where the CM field is always
$\Q(\sqrt{\Delta_{2}})$, in the cubic case the splitting field is sextic
generically and only quadratic-imaginary as a subfield in case~(1). The
``CM'' label therefore extends honestly only to case~(1)
($\Delta_{3}<0$, $S_{3}$ complex); cases~(2) and~(3) are totally real
and the analogous arithmetic obstruction, if it exists, must be
formulated via the cubic Hilbert class field of $L_{b}$ rather than via
imaginary quadratic CM theory. This is the most substantial structural
departure from PCF-1.
\end{remark}

\begin{conjecture}[B$_{3}$ part (iii): cubic Stokes signal]\label{conj:B3}
For a degree-three PCF in the irreducible-cubic-CM stratum (case~(1) of
Conjecture~\ref{conj:B2}: $\Delta_{3}<0$), there exists a one-parameter
deformation $K(t)$ that preserves $\Delta_{3}<0$ in a neighbourhood of
$t=0$ and respects the splitting-field structure, under which the Stokes
exponent
\[
S(t)\;=\;\frac{\log|L(t)-L(0)|}{\log|t|}
\]
satisfies $\lim_{t\to 0}S(t)\in(0,1)$. The conjectural fractional
exponent is governed by the local exponent differences of the
associated Heun-type cubic ODE.
\end{conjecture}

\begin{conjecture}[B$_{3}$ part (iv): Painlev\'e-hierarchy reduction]\label{conj:B4}
Each degree-three PCF family in the cell of Conjecture~\ref{conj:B3}
(irreducible $\Z$-primitive, $\Delta_{3}<0$, $\Gal(b/\Q)=S_{3}$ complex)
admits a Painlev\'e reduction in the P-II--P-VI hierarchy. Heuristic
alignment with cubic Galois data
(cf.~Iwasaki--Kimura--Shimomura--Yoshida~\cite{Iwasaki1991} and
Mazzocco~\cite{Mazzocco2003}) suggests P-IV as the most likely
instance, but no commitment to a specific Painlev\'e equation is made
at the program-statement level. The channel question is itself open:
PCF-1 Session~E has shown that the natural recurrence-parameter
deformation channel and the Borel-resummation channel both
\emph{fail} for the degree-2 $\Delta<0$ families with respect to
V$\_$quad's known P-III$(D_{6})$ reduction; the cubic case is
therefore expected to require a new channel construction, not a
verbatim lift of either PCF-1 channel.
\end{conjecture}

\begin{remark}\label{rem:painleveIV}
The heuristic preference for P-IV (over P-V or P-VI) within the
hierarchy rests on three structural observations, each of which is
necessary but not sufficient: (i) P-IV's isomonodromic problem has a
single irregular singularity of Poincar\'e rank $\frac{3}{2}$, the
minimal rank compatible with cubic-degree growth of $b_{n}$;
(ii) the B\"acklund-transformation group of P-IV is $A_{2}^{(1)}$, the
affine analogue of the cubic Galois group $A_{3}=C_{3}$;
(iii) P-IV admits Hermite/parabolic-cylinder special-function solutions
at half-integer parameters, the natural reservoir of explicit content
for the totally-real cyclic case~(B) of Conjecture~\ref{conj:B1}. None
of (i)--(iii) forces P-IV uniquely; the cubic case may equally well
fall into a P-V degeneration or a P-VI specialization at an
arithmetic point.
\end{remark}

\section{Worked Landmarks: Ap\'ery, Catalan-Cubic, Cyclic-$C_{3}$}\label{sec:landmarks}

To make the trichotomy of Conjecture~\ref{conj:B1} concrete we record
three landmark families that anchor the three Galois-group cells. None
is an empirical result of PCF-2; each is a known PCF whose role under
the proposed predicate is fixed in advance, so that Sessions A--C can be
calibrated against them.

\begin{table}[h]
\centering
\small
\begin{tabular}{p{0.13\textwidth}p{0.21\textwidth}rp{0.18\textwidth}p{0.10\textwidth}p{0.20\textwidth}}
\toprule
Landmark & $b(n)$ & $\Delta_{3}$ & Reducible? & $\Gal(b/\Q)$ & Predicted cell \\
\midrule
Ap\'ery $\zeta(3)$ & $34n^{3}+51n^{2}+27n+5$ & $-459=-3^{3}\cdot 17$ & yes: $(2n+1)(17n^{2}+17n+5)$ & $C_{2}$ on $\Q(\sqrt{-51})$ & out-of-scope (see Rmk.~\ref{rem:apery-reducible}) \\
Catalan-cubic seed & $n^{3}+n+1$ & $-31$ & no & $S_{3}$ (complex) & A: PSLQ-null + Stokes \\
Cyclic-$C_{3}$ probe & $n^{3}-3n+1$ & $81=9^{2}$ & no & $C_{3}$ & B: closed-form candidate \\
\bottomrule
\end{tabular}
\caption{Three landmark $b(n)$. The Ap\'ery row uses the recurrence
denominator $b(n)=34n^{3}+51n^{2}+27n+5$ from
\cite{vanderPoorten1979,Apery1979,Beukers1979}; an exact symbolic
computation (sympy, exact integer arithmetic) gives
$\Delta_{3}=\disc(b)=-459=-3^{3}\cdot 17$, and the polynomial
\emph{factors} over $\Q$ as $(2n+1)(17n^{2}+17n+5)$. Hence $b$ is
\emph{reducible}, the splitting field is
$\Q(\sqrt{17^{2}-4\cdot 17\cdot 5})=\Q(\sqrt{-51})$, and the Galois
group of $b$ over $\Q$ is $C_{2}$ rather than $C_{3}$ or $S_{3}$. The
Catalan-cubic seed $b(n)=n^{3}+n+1$ is the simplest irreducible cubic
with negative discriminant ($\disc=-31$, $\Gal=S_{3}$, splitting field
$\Q(\theta,\sqrt{-31})$); it is the cleanest test of
Conjecture~\ref{conj:B1}(A). The cyclic-$C_{3}$ probe $b(n)=n^{3}-3n+1$
is the standard simplest totally-real cyclic cubic with
$\disc=81=9^{2}$; it is the cleanest test of
Conjecture~\ref{conj:B1}(B). The accompanying $a_{n}$ for each landmark
is fixed in Session~A by the smallest $\Z$-primitive choice consistent
with the balanced-class convention~\eqref{eq:an}.}
\label{tab:landmarks}
\end{table}

\paragraph{Calibration role.}
Table~\ref{tab:landmarks} is a calibration anchor, not a results table:
each row records what the predicate \emph{must} say if Conjecture~\ref{conj:B1}
is to be consistent with literature. The Catalan-cubic and cyclic-$C_{3}$
rows are the positive-test equivalents for B3(i) cells (A) and (B); the
Ap\'ery row, despite its historical centrality at $d=3$, is
\emph{out of scope} of B3(i) by reducibility and thus contributes a
separate calibration role discussed in Remark~\ref{rem:apery-reducible}.

\begin{remark}[Ap\'ery is reducible, hence outside
B3(i)]\label{rem:apery-reducible}
The Ap\'ery denominator $b(n)=34n^{3}+51n^{2}+27n+5$ factors as
$(2n+1)(17n^{2}+17n+5)$, so it does not satisfy the irreducibility
hypothesis of Conjecture~\ref{conj:B1}. Ap\'ery's recurrence sits in
the \emph{out-of-scope, splittable-cubic} bin of trichotomy B3(i), and
its known identification $\lim K = 6/\zeta(3)$
\cite{vanderPoorten1979,Apery1979,Beukers1979} is consistent with the
restriction of Conjecture~\ref{conj:B1} to irreducible $\Z$-primitive
$b$: a reducible $b$ admits an algebraic factorization of the
recurrence and is not predicted to obstruct a closed form. The CM-style
subfield $\Q(\sqrt{-51})$ from the quadratic factor records the
arithmetic flavour but does not contradict the predicate. Operationally,
Session~A's calibration anchor is therefore the Ap\'ery family run
through the existing PCF-1 PSLQ pipeline as a \emph{positive control}
(target relation $L\sim\zeta(3)$ should be detected) for infrastructure
integrity, while the predicate B3(i) itself is tested on the
irreducible-only enumeration.
\end{remark}

\paragraph{Cubic Galois-group quick checks.}
For the irreducible $\Z$-primitive case, $\Gal(b/\Q)$ is cyclic ($C_{3}$)
iff $\disc(b)$ is a perfect square in $\Q$ (equivalently in $\Z$ for
$\Z$-primitive $b$); otherwise it is $S_{3}$. Examples:
\begin{itemize}[leftmargin=*]
\item $b(n)=n^{3}-3n+1$: $\disc=81=9^{2}$, $C_{3}$.
\item $b(n)=n^{3}+n+1$: $\disc=-31$, $S_{3}$ (complex case).
\item $b(n)=n^{3}-2$: $\disc=-108$, $S_{3}$ (complex case, but the
splitting field $\Q(\sqrt[3]{2},\zeta_{3})$ is the canonical $S_{3}$
cubic of cyclotomic origin).
\item $b(n)=n^{3}-n-1$: $\disc=-23$, $S_{3}$ (complex case, Lehmer's cubic).
\end{itemize}
The square-discriminant filter applied to $\Z$-primitive irreducible
cubics with $|\alpha_{i}|\le 5$ is heuristically rare (cyclic cubic
fields are sparse in the natural Davenport--Heilbronn / Bhargava
counts; see~\cite{Cohen2007} \S6.4 for the relevant background on
cubic discriminants and Galois groups), and the $\Delta_{3}<0$ fraction
is roughly $50\%$ by sign symmetry of the discriminant. The Session~A
enumeration is expected to populate cells~(A) (S$_{3}$ complex) and
(C) (S$_{3}$ totally real) of Conjecture~\ref{conj:B1} densely, with
cell~(B) (cyclic) populated sparsely; the calibration of these counts
against a finite-coefficient enumeration is itself part of the
Session~A deliverable.

\section{Computational Plan}\label{sec:plan}

The computational plan reuses the PCF-1 infrastructure: the Wallis
forward-recurrence routine, the PSLQ wrapper with the $18$-element basis
of PCF-1, the Stokes-exponent diagnostic from
PCF-1~\S 3, and the AEAL evidence-class scaffolding. The plan is
organized as three concrete sessions.

\paragraph{PCF-2 Session A: Cubic family enumeration and $\Delta_{3}$
classification.}
Enumerate $\Z$-primitive irreducible cubics $b(n)=\alpha_{3}n^{3}+\alpha_{2}n^{2}+\alpha_{1}n+\alpha_{0}$
with bounded coefficients $|\alpha_{i}|\le 5$ (target: $\sim 200$
families after gcd reduction and irreducibility filtering). For each
family, compute $\Delta_{3}$, classify the Galois group via the
$\sqrt{\Delta_{3}}\in\Q$ test, and tabulate the splitting type. Produce
a stratification table analogous to PCF-1
Table~1~\cite{Papanokechi2026PCF1}, with columns $\Delta_{3}$,
$\sqrt{\Delta_{3}}\in\Q?$, splitting type, Galois group, real-root count.
Estimated effort: 1 week.

\paragraph{PCF-2 Session B: PSLQ scan on irreducible-cubic-CM families.}
For each family in the $\Delta_{3}<0$ stratum produced by Session A
(expected: $\sim 80$ families), compute $L=K(a_{n},b_{n})$ to
$\ge 320$ stable digits via the Wallis recurrence at depth $2500$, and
run PSLQ against the PCF-1 $18$-element basis at working precision
$220$ digits and integer bound $10^{10}$. Confirm or refute
Conjecture~\ref{conj:B1}(A). For the $\sqrt{\Delta_{3}}\in\Q$ subset
of $\Delta_{3}>0$ families ($C_{3}$ totally-real cyclic),
run an extended PSLQ with the basis augmented by cubic-period
candidates (Catalan-style $\Gamma$-quotients at $1/3,2/3$;
$\sqrt[3]{2}$; $\sqrt[3]{3}$; ${}_{2}F_{1}$ values at $1/3$;
$\zeta(3)$ already in the base basis) at integer bound $10^{12}$ and
working precision $\ge 320$ digits. The targeted detection rate for
$C_{3}$ families is bounded above by the literature on cubic-period
identifications and may be substantially below~$1$. Estimated effort:
2 weeks.

\paragraph{PCF-2 Session C: Stokes pipeline for cubic case.}
For 6--10 families flagged in Session B as PSLQ-non-detected, design a
Galois-respecting one-parameter deformation, compute $L(t)$ at 11
$t$-values with the same precision/depth budget as the PCF-1 QL01
analysis, and compute $S(t)$ and $\sigma(t)$ via the formulas of PCF-1
\S 3. Confirm or refute Conjecture~\ref{conj:B3}. For the
strongest Stokes-signal candidate, attempt a P-IV ansatz fit at the
same overdetermined-design level used for V$\_$quad in PCF-1, and
record the residual matrix. Estimated effort: 2--3 weeks.

\paragraph{Precision and depth budget.}
The PCF-1 budget (working precision $220$ for PSLQ, $320$ for the
limit; depth $2500$; integer bound $10^{10}$) extends without
modification to $d=3$ for the PSLQ stage, since the basis size and the
expected coefficient bound do not change. The limit-computation budget
should be scrutinized at the outset of Session~A: cubic $b_{n}$
amplifies floating-point error growth in the Wallis recurrence by an
extra factor of $n$, so the depth-$2500$ stable-digit count may drop
from $\ge 320$ (PCF-1) to $\sim 280$ at $d=3$ unless precision is
raised correspondingly. We allocate working precision $360$ (limit) and
$280$ (PSLQ) as the conservative defaults, and treat the PCF-1
defaults as fallback.

\paragraph{Risk register.}
Three risks are identified in advance and each is matched to a
mitigation under the AEAL halt-conditions protocol.

\begin{table}[h]
\centering
\small
\begin{tabular}{p{0.30\textwidth}p{0.30\textwidth}p{0.32\textwidth}}
\toprule
Risk & Likelihood & Mitigation \\
\midrule
Wallis recurrence loses precision faster than expected at $d=3$, so
PSLQ-null results are not actually null but precision-starved. & Medium. & Run Ap\'ery landmark first; if $\zeta(3)/6$ is not detected at the conservative budget, escalate dps to $500$ and depth to $4000$. \\
Cubic-period basis is too small: $C_{3}$-cyclic limits identify against
constants outside the augmented basis. & Medium--high. & Treat all $C_{3}$ ``no relation'' as pre\-li\-mi\-na\-ry; flag for a follow-up cubic-Hilbert-class-field basis sweep before claim promotion. \\
P-IV ansatz fits all fail in Session~C, mirroring the QL01 outcome of
PCF-1. & High (this is what happened in the $d=2$ case). & Fall back to the Stokes-exponent $S<1$ predicate alone, weakening Conjecture~\ref{conj:B4} to a fractional-power existence claim, exactly as PCF-1 v1.2 did for QL01--QL26. \\
\bottomrule
\end{tabular}
\caption{Risk register for the PCF-2 computational plan.}
\label{tab:risks}
\end{table}

\paragraph{AEAL evidence chain.}
The session chain follows the PCF-1 pattern. Session~A is a
\emph{formalized} layer (the Galois-group classification is algorithmic
and exact). Session~B is a \emph{near\_miss} layer (each PSLQ
non-detection at sufficient precision is a near-miss-class evidence
artefact). Session~C is a \emph{numerical\_identity} layer (the Stokes
exponent and any P-IV parameter fit are numerical-identity-class
artefacts). \emph{Independently\_verified} promotion requires
cross-session reproduction of $L(0)$ to $\ge 50$ digits across at least
two independent sessions, identical to the PCF-1 cross-session protocol.

\section{Why This Matters}\label{sec:why}

PCF-2 is positioned to deliver one of two outcomes, both informative.

\paragraph{Outcome U (universality).}
If Conjecture~\ref{conj:B1} is confirmed across the Session~B sweep,
the $\Delta$-discriminant transcendence predicate of PCF-1 lifts cleanly
to higher degree, with the cubic Galois-group axis as the natural
refinement. This would be the first non-trivial empirical evidence that
the discriminant predicate is degree-general rather than a degree-two
coincidence, and it would provide a roadmap for PCF-3 ($d=4$, with the
quartic Galois groups $S_{4}, A_{4}, D_{4}, V_{4}, C_{4}$) and beyond.
The universality outcome also strengthens the analytic prediction of
Conjecture~\ref{conj:B4}: a Painlev\'e equation per degree, indexed by
the Galois group, would be a striking new entry in the
PCF--Painlev\'e correspondence catalogued by~\cite{Mazzocco2003}.

\paragraph{Outcome S (specificity).}
If Conjecture~\ref{conj:B1}(A) fails --- in particular, if any
$\Delta_{3}<0$ family yields a non-trivial PSLQ relation against the
PCF-1 basis --- then degree two is identified as a special boundary at
which the discriminant predicate works, and the higher-degree picture is
controlled by a different invariant (Picard--Fuchs operator structure,
modular-monodromy condition, or a cubic-period basis that supersedes
the PCF-1 $18$-element basis). The Ap\'ery family
($a_{n}=n^{6}$, $b_{n}=34n^{3}-51n^{2}+27n-5$) is the historical exemplar
of a $d=3$ PCF whose limit $\zeta(3)/6$ is detected against a cubic-period
basis; PCF-2 would then be repositioned as a search for the natural
generalization of the Ap\'ery basis. This outcome is no less valuable: it
sharpens the conjectural content of PCF-1 by identifying the precise
scope of the discriminant predicate.

\paragraph{Either way, a stable scope-anchor.}
The role of this program statement is to fix the scope, the conventions,
and the success criteria \emph{before} the computational sweep begins.
The AEAL methodology requires that conjectures be stated explicitly and
that evidence-class promotion follow a documented chain; this document
discharges the first half of that requirement. Sessions A--C of
\S\ref{sec:plan} discharge the second half.

\section{Relation to Existing Literature}\label{sec:litrel}

The mathematical content of PCF-2 sits at the intersection of three
literatures, none of which addresses the cubic PCF transcendence
question directly.

\paragraph{Painlev\'e and number theory.}
Mazzocco's lecture notes on Painlev\'e equations and number theory
\cite{Mazzocco2003} catalogue the algebraic-solution loci of P-VI
(and by reduction, P-V, P-IV, P-III) and identify a small set of
arithmetic Painlev\'e parameters at which classical special-function
solutions exist. The PCF-1 V$\_$quad result lifts naturally into
this framework as a P-III$(D_{6})$ solution at a specific
$\Q(\sqrt{-11})$-arithmetic point. The cubic analogue
(Conjecture~\ref{conj:B4}) predicts that P-IV --- whose isomonodromic
problem has Poincar\'e rank $3/2$ at the irregular singularity, the
natural cubic refinement of P-III's rank $1$ --- governs the cubic
case. Mazzocco's Backlund-transformation analysis of P-IV at $A_{2}^{(1)}$
parameters is the structural object PCF-2 Session~C is designed to
intersect.

\paragraph{Iwasaki--Kimura--Shimomura--Yoshida and the Gauss-to-Painlev\'e
arc.}
The monograph~\cite{Iwasaki1991} traces the lineage from the
Gauss hypergeometric equation to the Painlev\'e transcendents via
isomonodromic deformation. Each rank-two Heun reduction of a Painlev\'e
equation has a dedicated treatment, and the cubic-Galois-group
classification of Schwarz-triangle / Takeuchi groups
appears explicitly. The PCF-1 framework borrows the Spec$(K)$
$\tau_{\mathrm{arith}}$ component from this lineage; PCF-2 extends it to
the cubic Galois-group axis.

\paragraph{Ap\'ery's irrationality proof and the cubic PCF
exemplar.}
Van der Poorten's account~\cite{vanderPoorten1979} of
Ap\'ery's proof~\cite{Apery1979} of the irrationality of $\zeta(3)$ is
the historical landmark of cubic PCFs. The Ap\'ery recurrence is a
$d=3$ PCF whose limit is $6/\zeta(3)$, identified against a basis
that contains $\zeta(3)$ explicitly; this is the unique $d=3$ family
in the literature that is fully resolved against a closed-form basis
known in advance. Beukers~\cite{Beukers1979} reformulated the proof
via a Picard--Fuchs operator on a one-parameter family of K3
surfaces. The Picard--Fuchs framing is the candidate alternative
predicate raised under outcome~S of \S\ref{sec:why}: if the
discriminant predicate fails at $d=3$, the natural competitor is
``$L$ is identifiable iff the associated Picard--Fuchs operator
is a pullback of a hypergeometric operator over $\Q$''. PCF-2
Session~B is the test that decides between these two predicates.

\paragraph{Ramanujan-Machine search and PSLQ infrastructure.}
The Ramanujan-Machine programme~\cite{Raayoni2021} produced the bulk of
the PCF candidate database used by PCF-1, and its PSLQ-based
identification pipeline is what PCF-1 inherits. The
Ferguson--Bailey PSLQ algorithm~\cite{FergusonBailey1992} is
basis-dependent: a closed-form identification is only as strong as the
basis it is checked against. The cubic-period augmentation of the
PCF-1 basis (\S\ref{sec:plan}, Session~B) is therefore not optional
but mandatory for the $C_{3}$ cell of Conjecture~\ref{conj:B1}; the
$\Gamma(1/3)$, $\Gamma(2/3)$, $\sqrt[3]{2,3}$, ${}_{2}F_{1}(\cdot;1/3)$
augmentation is the smallest basis under which any reasonable
$C_{3}$-cyclic limit could conceivably be detected.

\paragraph{Cohen, cubic fields, and the heuristic distribution.}
Cohen's textbook treatment of cubic fields and their discriminants
\cite{Cohen2007} provides the standard background for the heuristic
distribution of $\Gal(b/\Q)$ across $\Z$-primitive irreducible cubics:
cyclic cubics are sparse (Davenport--Heilbronn / Bhargava counts),
which is consistent with the $\sim 5$--$10\%$ figure used in
\S\ref{sec:landmarks} as a Session~A pre-estimate. PCF-2 does not
rely on a published PCF-specific cubic enumeration; the Session~A
enumeration of \S\ref{sec:plan} performs its own bounded-coefficient
sweep and reports the empirical distribution as a deliverable.

\section{Acknowledgements}\label{sec:ack}

This program statement is part of the SIARC research program. The author
thanks the AEAL/SIARC multi-agent methodology of~\cite{P8AEAL} for
providing the evidence-gate scaffolding under which both PCF-1 and the
present PCF-2 program statement were prepared. AI tools (GitHub Copilot
powered by Anthropic Claude Opus 4.7, and Anthropic Claude) were used
for literature synthesis, prose drafting, and consistency-checking
across the SIARC preprint stack; all conjecture formulations and editorial
decisions are the author's. The cubic-family enumeration plan of
Session~A reuses the \texttt{f1\_base\_computation.py} infrastructure of
PCF-1, and the Stokes-exponent diagnostic of Session~C reuses the
deformation-design protocol of PCF-1 \S 3. The Painlev\'e--cubic
literature consulted in formulating Conjecture~\ref{conj:B4} is
indicated in the references.

\begin{thebibliography}{99}

\bibitem{Papanokechi2026PCF1}
Papanokechi, \emph{Complex Multiplication as a Transcendence Predicate
for Degree-2 Polynomial Continued Fractions} (P12 / PCF-1, v1.2).
Zenodo preprint, 2026.
\href{https://doi.org/10.5281/zenodo.19934553}
{DOI:10.5281/zenodo.19934553}.

\bibitem{Papanokechi2026UMB}
Papanokechi, \emph{An Arithmetic Stratification of Polynomial Continued
Fractions: Program Statement of the SIARC Series}.
Zenodo preprint, 2026.
\href{https://doi.org/10.5281/zenodo.19885550}
{DOI:10.5281/zenodo.19885550}.

\bibitem{Papanokechi2026Spec}
Papanokechi, \emph{Spectral Classes of Polynomial Continued Fractions:
A Hypergeometric Framework and Arithmetic Obstructions}.
Preprint, 2026.

\bibitem{Papanokechi2026Vquad}
Papanokechi, \emph{Resurgent Analysis of the V$\_$quad Continued
Fraction}.
Zenodo preprints, 2026.
\href{https://doi.org/10.5281/zenodo.19491767}
{DOI:10.5281/zenodo.19491767} and
\href{https://doi.org/10.5281/zenodo.19491768}
{DOI:10.5281/zenodo.19491768}.

\bibitem{P7AIBehSci}
Papanokechi, \emph{AI Behavior Science: A Proposal for an Empirical
Discipline} (P7). Zenodo, 2026.
\href{https://doi.org/10.5281/zenodo.19562751}
{DOI:10.5281/zenodo.19562751}.

\bibitem{P8AEAL}
Papanokechi, \emph{AEAL: AI-Augmented Evidence Aggregation and
Logging --- A Multi-Agent Methodology for Empirical Mathematics} (P8).
Zenodo, 2026.
\href{https://doi.org/10.5281/zenodo.19565086}
{DOI:10.5281/zenodo.19565086}.

\bibitem{Apery1979}
R.~Ap\'ery, \emph{Irrationalit\'e de $\zeta(2)$ et $\zeta(3)$}.
Ast\'erisque~61 (1979), 11--13.

\bibitem{Beukers1979}
F.~Beukers, \emph{A note on the irrationality of $\zeta(2)$ and
$\zeta(3)$}. Bull.~London Math.~Soc.~11 (1979), 268--272.

\bibitem{Mazzocco2003}
M.~Mazzocco, \emph{Painlev\'e sixth equation as isomonodromic deformations
equation of an irregular system}, in: \emph{Isomonodromic Deformations and
Applications in Physics} (J.~Harnad, A.~R.~Its, eds.),
CRM Proc.~Lecture Notes~31, Amer.~Math.~Soc., 2003, 219--238.
See also: M.~Mazzocco, \emph{Painlev\'e equations and number theory},
lecture notes.

\bibitem{Iwasaki1991}
K.~Iwasaki, H.~Kimura, S.~Shimomura, M.~Yoshida,
\emph{From Gauss to Painlev\'e: A Modern Theory of Special Functions}.
Aspects of Mathematics E16, Vieweg, Braunschweig, 1991.

\bibitem{Okamoto1987}
K.~Okamoto, \emph{Studies on the Painlev\'e equations IV: Third
Painlev\'e equation $P_{III}$}. Funkcial.~Ekvac.~30 (1987), 305--332.

\bibitem{Raayoni2021}
G.~Raayoni et~al., \emph{Generating conjectures on fundamental constants
with the Ramanujan Machine}.
Nature~590 (2021), 67--73.

\bibitem{FergusonBailey1992}
H.~R.~P.~Ferguson and D.~H.~Bailey, \emph{A polynomial-time, numerically
stable integer relation algorithm}. RNR Tech.\ Rep.\ RNR-91-032, NASA
Ames, 1992.

\bibitem{vanderPoorten1979}
A.~van der Poorten, \emph{A proof that Euler missed\ldots\ Ap\'ery's proof
of the irrationality of $\zeta(3)$}. Math.~Intelligencer~1 (1978/79),
195--203.

\bibitem{Cohen2007}
H.~Cohen, \emph{Number Theory, Volume I: Tools and Diophantine
Equations}. Graduate Texts in Mathematics~239, Springer-Verlag, New
York, 2007.

\end{thebibliography}

\end{document}
